%% LyX 2.0.2 created this file.  For more info, see http://www.lyx.org/.
%% Do not edit unless you really know what you are doing.
\documentclass[english]{article}
\usepackage[T1]{fontenc}
\usepackage[latin9]{inputenc}
\usepackage{amsthm}
\usepackage{amsmath}
\usepackage{amssymb}
\PassOptionsToPackage{normalem}{ulem}
\usepackage{ulem}

\makeatletter
%%%%%%%%%%%%%%%%%%%%%%%%%%%%%% Textclass specific LaTeX commands.
  \theoremstyle{plain}
  \newtheorem{lem}{\protect\lemmaname}
  \theoremstyle{definition}
  \newtheorem{defn}{\protect\definitionname}
  \theoremstyle{plain}
  \newtheorem{cor}{\protect\corollaryname}
  \theoremstyle{remark}
  \newtheorem{rem}{\protect\remarkname}

\makeatother

\usepackage{babel}
  \providecommand{\definitionname}{Definition}
  \providecommand{\lemmaname}{Lemma}
  \providecommand{\remarkname}{Remark}
\providecommand{\corollaryname}{Corollary}

\begin{document}

\title{RTS Dev by The Formal Approach - Final Assignment}


\author{Yoni Zohar 066480641}

\maketitle

\section{Part A - Problem B}
\begin{lem}
\label{lem: the first disjunct}If $s_{0}\models\psi_{2}\wedge EX_{\left[a,b\right]}\left(true\right)$
then $s_{0}\models E\left(\psi_{1}U_{\left[a,b\right]}\psi_{2}\right)$.\end{lem}
\begin{proof}
Since $s_{0}\models EX_{\left[a,b\right]}\left(true\right)$, this
means that there exists $s_{1}$ such that $\left\langle s_{0},n,s_{1}\right\rangle \in R$
and $a\leq n\leq b$ (and $s_{1}\models true$, but this isn't important).
In particular, consider any path starting at $s_{1}$ and denote it
as $s_{1},s_{2},...$. Consider the path $s_{0},s_{1},s_{2},...$.
Let $k=0$. Indeed, for this $k$, $s_{k}\models\psi_{2}$, there
are no vertices between $s_{0}$ and $s_{k}$ and $a\leq\Sigma_{i=0}^{k}s_{i}=s_{0}\leq b$.
Hence $s_{0}\models E\left(\psi_{1}U_{\left[a,b\right]}\psi_{2}\right)$.\end{proof}
\begin{defn}
Let $\varphi$ be a $TTG-CTL$ formula and let $K$ be a $TTG$. $K\left(\varphi\right):=\left\{ s\in K:\, s\models\varphi\right\} $.\end{defn}
\begin{lem}
\label{lem: EU =00003D or and EU}Let $K=\left\langle S,s_{I},R,L\right\rangle $
be a $TTG$ and let $\psi_{1},\psi_{2}$ be $TTG-CTL$ formulas. Let
$a\leq b\in\mathbb{N}$. Then there are $c,d,e$ such that:
\begin{itemize}
\item $K\left(E\left(\psi_{1}U_{\left[a,b\right]}\psi_{2}\right)\right)=K\left(\left(\psi_{2}\wedge EX_{\left[a,b\right]}\left(true\right)\right)\vee\left(\psi_{1}\wedge EX_{\left[c,c\right]}E\left(\psi_{1}U_{\left[d,e\right]}\psi_{2}\right)\right)\right)$.
\item $a\leq c+d$
\item $c+e\leq b$.
\end{itemize}
\end{lem}
\begin{proof}
$\left(\Leftarrow\right)$ Suppose $s\models\left(\psi_{2}\wedge EX_{\left[a,b\right]}\left(true\right)\right)\vee\left(\psi_{1}\wedge EX_{\left[c,c\right]}E\left(\psi_{1}U_{\left[d,e\right]}\psi_{2}\right)\right)$
for some $c,d,e$ such that $a\leq c+d$ and $c+e\leq b$. If $s\models\left(\psi_{1}\wedge EX_{\left[a,b\right]}\left(true\right)\right)$
then by Lemma \ref{lem: the first disjunct} $s\models E\left(\psi_{1}U_{\left[a,b\right]}\psi_{2}\right)$.
Otherwise, \end{proof}
\begin{lem}
Let $K=\left\langle S,s_{I},R,L\right\rangle $ be a $TTG$ and let
$\psi_{1},\psi_{2}$ be $TTG-CTL$ formulas. Let $a\leq b\in\mathbb{N}$.
Then $K\left(A\left(\psi_{1}U_{\left[a,b\right]}\psi_{2}\right)\right)=K\left(\psi_{2}\vee\left(\psi_{1}\wedge AXA\left(\psi_{1}U_{\left[a,b\right]}\psi_{2}\right)\right)\right)$.\end{lem}
\begin{proof}
Similar to the proof of Lemma \ref{lem: EU =00003D or and EU}. \end{proof}
\begin{cor}
~
\begin{itemize}
\item $E\left(\psi_{1}U_{\left[a,b\right]}\psi_{2}\right)$ is a fix-pint
of $\lambda Z.\psi_{2}\vee\left(\psi_{1}\wedge EXZ\right)$.
\item \textup{$A\left(\psi_{1}U_{\left[a,b\right]}\psi_{2}\right)$ is a
fix-point of $\lambda Z.\psi_{2}\vee\left(\psi_{1}\wedge AXZ\right)$.}
\end{itemize}
\end{cor}
\begin{rem}
From now on we fix a $TTG$ $K=\left\langle S,S_{I},R,L\right\rangle $
and identify formulas with the set of states of that $TTG$ which
satisfy it. \end{rem}
\begin{lem}
The functional $\lambda Z.\psi_{2}\vee\left(\psi_{1}\wedge EXZ\right)$
is monotonic.\end{lem}
\begin{proof}
Assume $Z_{1}\subseteq Z_{2}$. We want to show that $\psi_{2}\vee\left(\psi_{1}\wedge EXZ_{1}\right)\subseteq\psi_{2}\vee\left(\psi_{1}\wedge EXZ_{1}\right)$.
First, we show that $EXZ_{1}\subseteq EXZ_{1}$. Assume $s\models EXZ_{1}$.
Then there exists $s'$ such that 
\end{proof}

\section{Part B}


\subsection{Problem Assumptions and Requirements}


\subsection{Ontology}
\begin{itemize}
\item $P\_sgnl$
\item $E\_sgnl$
\item $T\_cmd:=Red$
\item $T\_cmd:=Green$
\item $P\_cmd:=Red$
\item $P\_cmd:=Green$
\end{itemize}

\subsection{Abstract Variables}
\begin{itemize}
\item $T\_state:=\left(T\_cmd=Green\right)\wedge\left(P\_cmd=Red\right)$
\item $P\_state:=\left(T\_cmd=Red\right)\wedge\left(P\_cmd=Green\right)$ 
\item $I\_state:=\left(T\_cmd=Red\right)\wedge\left(P\_cmd=Red\right)$
\item $Bad\_state:=\left(T\_cmd=Green\right)\wedge\left(P\_cmd=Green\right)$
\end{itemize}
~
\begin{itemize}
\item $T\_enters:=\left(\left(\approx\left(\neg\left(T\_state\right)\right)\right)\wedge\left(T\_state\right)\right)\vee\left(\left(\approx\left(T\_state\right)\right)\wedge\left(E\_sgnl\right)\wedge\left(T\_state\right)\right)$
\item $P\_enters:=\left(\left(\approx\left(\neg P\_state\right)\right)\wedge P\_state\right)\vee\left(\left(\approx\left(P\_state\right)\right)\wedge\left(P\_sgnl\right)\wedge\left(P\_state\right)\right)$
\end{itemize}

\subsubsection{Assumptions}
\begin{itemize}
\item An internal mechanism of $P$-button disables $P$-signal after one
press in $P$-state. $P$-signal generation is re-enabled only upon
the entrance to $T$-state. 

\begin{description}
\item [{LTL}] $\Box\left(\left(\left(\left(P\_state\right)\vee\left(\approx\left(P\_state\right)\right)\right)\wedge\left(P\_sgnl\right)\right)\rightarrow\left(\left(\neg P\_sgnl\right)\, U\,\left(T\_enters\right)\right)\right)$
\footnote{The actual P-state is continued for (almost) one more moment. Therefore,
the $P\_sgnl$ count includes the (almost) one moment after the $P\_state$
(right before the transition to $I$-state).%
}
\item [{XCTL}] 
\end{description}
\end{itemize}

\subsubsection{Requirements }
\begin{itemize}
\item Transfer from $T$-state to $P$-state must be seperated by $I$-state.

\begin{description}
\item [{LTL}] $\Box\left(T\_state\rightarrow X\left(\neg\left(P\_state\right)\right)\right)\wedge\Box\left(\left(\left(T\_state\right)\wedge X\left(X\left(P\_state\right)\right)\right)\rightarrow\left(X\left(I\_state\right)\right)\right)$.
\item [{XCTL}] 
\end{description}
\item Transfer from $P$-state to $T$-state must be seperated by $I$-state.

\begin{description}
\item [{LTL}] $\Box\left(P\_state\rightarrow X\left(\neg\left(T\_state\right)\right)\right)\wedge\Box\left(\left(\left(P\_state\right)\wedge X\left(X\left(T\_state\right)\right)\right)\rightarrow X\left(I\_state\right)\right)$.
\item [{XCTL}] 
\end{description}
\item The system should never be in a state where $T$-light and $P$-light
show green, both: 

\begin{description}
\item [{LTL}] $\Box\left(\neg\left(Bad\_state\right)\right)$
\item [{XCTL}] 
\end{description}
\item Initially, the system must be at $T$-state.

\begin{description}
\item [{LTL}] $T\_state$
\item [{XCTL}] 
\end{description}
\item Each time the system enters $T$-state it should stay there for at
least 4 seconds.

\begin{description}
\item [{LTL}] $\Box\left(T\_enters\rightarrow\left(X^{0,3}\left(T\_state\right)\right)\right)$%
\footnote{In my formulation, $T\_cmd:=Red$ means that at this moment this command
is issued. This means that if this command is issued 3 seconds in
a row, the $T$ light will be red for 4 seconds (It will start to
be red now for one second and then restart 3 more times, each time
for one more second)%
} 
\item [{XCTL}] 
\end{description}
\item The system should never leave $T$-state for more than 7 seconds.

\begin{description}
\item [{LTL}] $\Box\left(\left(\left(\approx\left(T\_state\right)\right)\wedge\neg T\_state\right)\rightarrow\left(X_{1,7}\left(T\_state\right)\right)\right)$
\item [{XCTL}] 
\end{description}
\item The system shall transfer from $T$-state to $P$-state only in response
to $P$-signal.

\begin{description}
\item [{LTL}] $\Box\left(\neg\left(P\_sgnl\right)\right)\rightarrow\Box\left(\neg\left(P\_state\right)\right)$
\item [{XCTL}] 
\end{description}
\item The system should reach $P$-state within 5 seconds since the occurence
of $P$-signal.

\begin{description}
\item [{LTL}] $\Box\left(P\_sgnl\rightarrow\left(X_{0,5}\left(P\_state\right)\right)\right)$
\footnote{Since P\_state means that the lights command is issued now, we want
that the commands will be issued in the next 5 seconds. Therefore
we don't use the same approach as above. There, we diminished the
interval by one. Here, for the same reason as above, we don't diminish
it. In both cases the main point is that a command is ``true'' in
a state if it is started or restarted at that state. %
}
\item [{XCTL}] 
\end{description}
\item If an $E$-signal occurs while the system is in $I$-state - in transfer
from $T$-state to $P$-state then the system should return immediately
to $T$-state. 

\begin{description}
\item [{LTL}] $\Box\left(\left(\left(I\_state\right)\wedge\left(\approx\left(T\_state\right)\right)\wedge\left(E\_sgnl\right)\right)\rightarrow\left(X\left(T\_enters\right)\right)\right)$
\item [{XCTL}] 
\end{description}
\item Once the system enters $P$-state it should stay there 3 seconds,
unless there was an occurrence of an $E$-signal

\begin{description}
\item [{LTL}] $\Box\left(\left(P\_enters\right)\rightarrow\left(\left(X^{0,3}\left(\neg E\_sgnl\right)\right)\rightarrow\left(X^{0,3}\left(P\_state\right)\right)\right)\right)$ 
\item [{XCTL}] 
\end{description}
\item The stay in $P$-state shall be constrained however by an occurrence
of $E$-signal, in which case the system would reach $T$-state within
2 seconds, not necessarily respecting the requirement of 3 seconds
stay in $P$-state.

\begin{description}
\item [{LTL}] $\Box\left(\left(P\_state\wedge E\_sgnl\right)\rightarrow X_{1,2}\left(T\_enters\right)\right)$
\item [{XCTL}] 
\end{description}
\item A $P$-signal that occurs while being in $I$-state, in transfer from
$P$-state to $T$-state should be served, but only after entering
$T$-state.

\begin{description}
\item [{LTL}] $\Box\left(\left(Y\left(I\_state\right)\wedge P\_sgnl\wedge\left(\approx\left(P\_state\right)\right)\right)\rightarrow X\left(T\_enters\wedge\left(X_{0,4}\left(P\_state\right)\right)\right)\right)$.%
\footnote{``while being in $I$-state'': This means that we've entered $I$-state
(i.e. $I\_state$ \textbf{\uline{was}} a true statement), and from
then until now a $P\_sgnl$ occurred. This is why I use $Y\left(I\_state\right)$.%
}
\item [{XCTL}] 
\end{description}
\end{itemize}

\subsubsection{Additional Assumptions}

The above specification is not consistent: For example, if from the
beginning both $E\_sgnl$ and $P\_sgnl$ are on non-stop, then by
one of the requirements ($\Box\left(T\_enters\rightarrow\left(X^{0,4}\left(T\_state\right)\right)\right)$)
we must stay forever at $T\_state$ (each second we must stay there
for the next 4 seconds, because of the $E\_sgnl$). By another one
of the requirements, we must reach $P\_state$ within 5 seconds ($\Box\left(P\_sgnl\rightarrow X_{0,5}\left(P\_state\right)\right)$).
These requirements can't coincide. Therefore, I am adding the following
assumption:
\begin{itemize}
\item The $P$-button mechanism has a very advanced (and internal) $E$-sensor.
This sensor disables the $P$-button 7 seconds \textbf{before} an
$E\_sgnl$ occurs. This sensor re-enables the $P\_button$ 7 second
\textbf{after} the $E\_sgnl$ occurs. 

\begin{description}
\item [{LTL}] $\Box\left(\left(E\_sgnl\right)\rightarrow\left(\left(X^{0,7}\left(\neg\left(P\_sgnl\right)\right)\right)\wedge\left(Y_{0,7}\left(\neg\left(P\_sgnl\right)\right)\right)\right)\right)$.
\item [{XCTL}] 
\end{description}
\end{itemize}
I think that this is a reasonable assumption. First, I think that
no assumptions should be made regarding the $E\_sgnl$, as we can't
predict emergencies. It is more important for the system to handle
emergencies well than to take care of pedestrians starvation. Pedestrians
that try to activate the $P\_sgnl$ with no success should try again
after a few seconds. If (god forbids) an $E\_sgnl$ keeps reoccurring,
then the pedestrians should just make a detour and get to their destination
using another way. As for how exactly this magic mechanism works,
there could be several possibilities: first, it could be that the
$P$-button $E$-sensor is more powerful than the $E$-sensor of the
system. Moreover, the $P$-button $E$-sensor might operate on sounds,
hearing the siren of the emergency car 5 seconds before it arrives
at the system's $E$-sensor. The extra 5 seconds are quite arbitrary,
and make the design simpler. It is possible to shorten the disablement
period. 


\subsection{NuSMV}


\subsection{State-Chart}
\end{document}
